\scanpage{20}
est représentable. En effet, laissant d'abord
tomber la dernière condition, on est
ramené à un sous-schéma \deleted{fermé} de
$\uuline{\mathrm{Hilb}}_{P/S}$ [\inserted{fermé} pour les sous-schémas de $P$
qui sont des sous-groupes de $\mathrm{Pic}$, et
ouvert dans ce dernier pour la condition
abélienne], montrons que la condition
ensembliste est une condition\unsure{,|\,} \uline{ouverte}:
elle est en effet invariante par changement
de base, et si elle est
vérifiée pour un $A\subset P$ en un pt,
elle l'est pour raison de dimensions,
dans un voisinage. De plus,
la projection $M\to S$ de ce problème
modulaire a les propriétés suivantes

(i) les fibres géométriques sont isomorphes \unsure{aux|à un|à}
\deleted{$\unsure{k(x)|k(y)}$} $\mathrm{Spec}\ k(x)$. Évident, en regardant
sur un corps alg. clos, et utilisant \unsure{les|le}
\uuline{th. de rigidité}

(ii) La projection $M\to S$ est propre.

Elle est en effet de type fini\unsure{;|,} séparée, reste
à appliquer le critère valuatif, savoir
ceci: dans le cas où $S$ est le spectre d'un
% archive sheet number: 19
